\chapter[Inner Product Spaces]{Inner Product Spaces and Spectral Theory}\label{chap:inner-products-spectral-theory}

Chapters~6 and~7 develop linear maps algebraically: kernels, images,
eigenvalues, and canonical forms do not require a notion of length. An inner
product adds geometry. It lets us ask stronger questions: when is a direct
sum orthogonal, when is a matrix a genuine change of orthonormal coordinates,
and when can an eigenbasis be chosen orthonormally? The answers lead to the
finite-dimensional spectral theorem and singular value decomposition.

Throughout this chapter \(V\) and \(W\) are finite-dimensional vector spaces
over \(F=\mathbb R\) or \(F=\mathbb C\). In the complex case our convention is
that the inner product is linear in the first variable and conjugate-linear in
the second. This convention agrees with the one used for adjoints below.

\section{Inner Products and Norms}

Length and angle in \(\mathbb R^n\) are encoded by the dot product. The
following axioms retain its useful properties while incorporating the complex
conjugation that makes a complex squared length real and nonnegative.

\begin{definition}
An \emph{inner product} on \(V\) is a map
\[
\langle\ ,\ \rangle:V\times V\longrightarrow F
\]
such that, for all \(u,v,w\in V\) and \(a,b\in F\),
\[
\langle au+bv,w\rangle
=a\langle u,w\rangle+b\langle v,w\rangle,
\]
\[
\langle v,u\rangle=\overline{\langle u,v\rangle},
\]
and
\[
\langle v,v\rangle\geq0,
\qquad
\langle v,v\rangle=0\Longleftrightarrow v=0.
\]
The associated \emph{norm} is
\[
\lVert v\rVert=\sqrt{\langle v,v\rangle}.
\]
\end{definition}

Conjugate symmetry turns first-variable linearity into
\[
\langle u,av+bw\rangle
=\overline a\langle u,v\rangle+\overline b\langle u,w\rangle.
\]
In particular,
\[
\lVert av\rVert=|a|\lVert v\rVert.
\]
Thus the complex convention is not cosmetic: it is what makes a scalar act
on length by its ordinary absolute value.

\begin{example}[Standard inner products]
On \(F^n\), the standard inner product is
\[
\langle u,v\rangle=\sum_{i=1}^n u_i\overline{v_i}.
\]
On \(M_{m,n}(F)\), the formula
\[
\langle A,B\rangle=\operatorname{tr}(AB^*)
=\sum_{i=1}^m\sum_{j=1}^n a_{ij}\overline{b_{ij}},
\qquad B^*=\overline B^{\,t},
\]
is an inner product. On the vector space \(P_n(F)\) of polynomials of degree
at most \(n\),
\[
\langle p,q\rangle=\int_0^1p(t)\overline{q(t)}\,dt
\]
is also an inner product. Positivity in the last example follows because a
nonzero polynomial is nonzero on some interval, so its squared absolute value
has strictly positive integral.
\end{example}

\begin{example}[Why conjugation is necessary]
The formula
\[
\sum_{i=1}^n u_iv_i
\]
is not an inner product on \(\mathbb C^n\): for \(u=(1,i)\), it gives
\[
1+i^2=0
\]
although \(u\ne0\). The standard formula instead gives
\[
\langle u,u\rangle=|1|^2+|i|^2=2.
\]
Conjugation is exactly what turns the complex dot product into a genuine
squared length.
\end{example}

\begin{theorem}[Cauchy--Schwarz inequality]
For all \(u,v\in V\),
\[
|\langle u,v\rangle|\leq\lVert u\rVert\lVert v\rVert.
\]
Equality holds, for nonzero \(u\) and \(v\), if and only if they are linearly
dependent.
\end{theorem}
\begin{proof}
If \(v=0\), the inequality is immediate. Otherwise put
\[
c=\frac{\langle u,v\rangle}{\lVert v\rVert^2}.
\]
Positivity gives
\[
0\leq\lVert u-cv\rVert^2.
\]
Using the linearity convention in the two variables, the right-hand side
expands to
\[
\lVert u\rVert^2
-\frac{|\langle u,v\rangle|^2}{\lVert v\rVert^2}.
\]
Multiplication by \(\lVert v\rVert^2\) proves the inequality. Equality
holds exactly when \(u-cv=0\), which is exactly linear dependence when both
vectors are nonzero.
\end{proof}

\begin{proposition}[Elementary norm identities]
For \(u,v\in V\),
\[
\lVert u+v\rVert\leq\lVert u\rVert+\lVert v\rVert,
\]
\[
u\perp v
\quad\Longrightarrow\quad
\lVert u+v\rVert^2=\lVert u\rVert^2+\lVert v\rVert^2,
\]
and
\[
\lVert u+v\rVert^2+\lVert u-v\rVert^2
=2\lVert u\rVert^2+2\lVert v\rVert^2.
\]
The last identity is the \emph{parallelogram identity}.
\end{proposition}
\begin{proof}
Expanding gives
\[
\lVert u+v\rVert^2
=\lVert u\rVert^2+2\operatorname{Re}\langle u,v\rangle+\lVert v\rVert^2.
\]
Cauchy--Schwarz bounds this by
\[
(\lVert u\rVert+\lVert v\rVert)^2,
\]
which proves the triangle inequality after taking nonnegative square roots.
When \(u\perp v\), the middle term vanishes, giving Pythagoras. Adding the
corresponding expansions for \(u+v\) and \(u-v\) cancels the two cross terms
and proves the parallelogram identity.
\end{proof}

\begin{proposition}[Polarization]
The norm induced by an inner product determines the inner product. More
precisely, over \(\mathbb R\),
\[
\langle u,v\rangle
=\frac14\bigl(\lVert u+v\rVert^2-\lVert u-v\rVert^2\bigr),
\]
whereas over \(\mathbb C\), with our first-variable-linear convention,
\[
\langle u,v\rangle
=\frac14\bigl(\lVert u+v\rVert^2-\lVert u-v\rVert^2\bigr)
+\frac{i}{4}\bigl(\lVert u+iv\rVert^2-\lVert u-iv\rVert^2\bigr).
\]
\end{proposition}
\begin{proof}
The first difference is
\[
\lVert u+v\rVert^2-\lVert u-v\rVert^2
=4\operatorname{Re}\langle u,v\rangle.
\]
Over \(\mathbb C\), direct expansion of the second difference gives
\[
\lVert u+iv\rVert^2-\lVert u-iv\rVert^2
=4\operatorname{Im}\langle u,v\rangle.
\]
The displayed formulas follow by combining real and imaginary parts. The
sign of the second term reflects the convention that the first, rather than
the second, variable is linear.
\end{proof}

\section{Orthonormal Sets and Gram--Schmidt}

Orthogonality makes coordinates behave like independent directions: all cross
terms disappear. Gram--Schmidt shows that this convenience is available in
every finite-dimensional inner-product space.

\begin{definition}
Vectors \(u,v\in V\) are \emph{orthogonal}, written \(u\perp v\), if
\[
\langle u,v\rangle=0.
\]
A set is \emph{orthogonal} if distinct members are orthogonal and
\emph{orthonormal} if, in addition, every member has norm one. For a subset
\(S\subseteq V\), its \emph{orthogonal complement} is
\[
S^\perp=\{v\in V:\langle v,s\rangle=0\text{ for every }s\in S\}.
\]
\end{definition}

\begin{proposition}[Orthogonal coordinates]\leavevmode\\
Every orthogonal set of nonzero vectors is independent. If
\[
(e_1,\ldots,e_n)
\]
is an orthonormal basis of \(V\), then every \(v\in V\) has the expansion
\[
v=\sum_{j=1}^n\langle v,e_j\rangle e_j,
\]
and
\[
\lVert v\rVert^2=\sum_{j=1}^n|\langle v,e_j\rangle|^2.
\]
\end{proposition}
\begin{proof}
If
\[
\sum_{j=1}^na_je_j=0
\]
for an orthogonal set of nonzero vectors, take the inner product with
\(e_k\). This gives
\[
a_k\lVert e_k\rVert^2=0,
\]
so every \(a_k=0\). For an orthonormal basis, write
\[
v=\sum_{j=1}^na_je_j.
\]
Taking the inner product with \(e_k\) gives \(a_k=\langle v,e_k\rangle\).
Pythagoras applied to this orthogonal sum gives the norm formula.
\end{proof}

\begin{corollary}[Bessel and Parseval]
If \(e_1,\ldots,e_r\) is an orthonormal set in \(V\), then
\[
\sum_{j=1}^r|\langle v,e_j\rangle|^2\leq\lVert v\rVert^2
\qquad(v\in V).
\]
If the set is an orthonormal basis, equality holds; this equality is
\emph{Parseval's identity} in the finite-dimensional setting.
\end{corollary}
\begin{proof}
Put
\[
w=\sum_{j=1}^r\langle v,e_j\rangle e_j.
\]
Then \(v-w\) is orthogonal to every \(e_j\), hence to \(w\). Pythagoras
gives
\[
\lVert v\rVert^2=\lVert w\rVert^2+\lVert v-w\rVert^2
=\sum_{j=1}^r|\langle v,e_j\rangle|^2+\lVert v-w\rVert^2.
\]
This proves Bessel's inequality. If the \(e_j\)'s form a basis, then
\(w=v\), so equality holds.
\end{proof}

\begin{theorem}[Gram--Schmidt orthogonalization]
Let
\[
(v_1,\ldots,v_n)
\]
be a linearly independent list. Define first
\[
u_1=v_1,
\qquad
e_1=\frac{u_1}{\lVert u_1\rVert},
\]
then
\[
u_2=v_2-\langle v_2,e_1\rangle e_1,
\qquad
e_2=\frac{u_2}{\lVert u_2\rVert},
\]
then
\[
u_3=v_3-\langle v_3,e_1\rangle e_1
-\langle v_3,e_2\rangle e_2,
\qquad
e_3=\frac{u_3}{\lVert u_3\rVert}.
\]
In general,
\[
u_k=v_k-\sum_{j=1}^{k-1}\langle v_k,e_j\rangle e_j,
\qquad
e_k=\frac{u_k}{\lVert u_k\rVert}.
\]
Every \(u_k\) is nonzero, the resulting list is orthonormal, and
\[
\operatorname{span}(e_1,\ldots,e_k)
=\operatorname{span}(v_1,\ldots,v_k)
\qquad(1\leq k\leq n).
\]
\end{theorem}
\begin{proof}
Assume that \(e_1,\ldots,e_{k-1}\) have been constructed. For \(i<k\),
\[
\langle u_k,e_i\rangle
=\langle v_k,e_i\rangle
-\sum_{j=1}^{k-1}\langle v_k,e_j\rangle\langle e_j,e_i\rangle
=0.
\]
Thus \(u_k\) is orthogonal to the preceding \(e_i\)'s. If \(u_k=0\), the
definition would put \(v_k\) in the span of \(e_1,\ldots,e_{k-1}\), which,
by the induction hypothesis, is the span of \(v_1,\ldots,v_{k-1}\). This
contradicts independence. Hence \(e_k\) is defined and has norm one.

The formula defining \(u_k\) shows that \(u_k\) lies in the span of the first
\(k\) vectors \(v_i\). Conversely, it can be rearranged to express \(v_k\)
in the span of \(u_k,e_1,\ldots,e_{k-1}\). The two successive spans are
therefore equal. Induction proves every assertion.
\end{proof}

\begin{corollary}[Gram--Schmidt fixes orthogonal lists]
If \(v_1,\ldots,v_k\) is an orthogonal list of nonzero vectors, then
Gram--Schmidt applied to this ordered list gives
\[
e_i=\frac{v_i}{\lVert v_i\rVert}
\qquad(1\leq i\leq k).
\]
In particular, Gram--Schmidt fixes every ordered orthonormal set.
\end{corollary}
\begin{proof}
At the \(i\)-th step, the preceding vectors \(e_j\) are scalar multiples of
\(v_j\).  Orthogonality gives
\[
\langle v_i,e_j\rangle=0
\qquad(j<i),
\]
so every projection term in the Gram--Schmidt formula vanishes.  Thus
\(u_i=v_i\), and normalizing gives the displayed formula.  If the list is
already orthonormal, then \(\lVert v_i\rVert=1\), so \(e_i=v_i\).
\end{proof}

\begin{corollary}[Orthonormal bases]\leavevmode\\
Each finite-dimensional inner-product space admits an orthonormal basis.
Every orthonormal set can be enlarged to one.
\end{corollary}
\begin{proof}
Apply Gram--Schmidt to any basis. For an orthonormal set, first extend it to
an ordinary basis using the finite-dimensional basis theorem of Chapter~5,
then apply Gram--Schmidt to the additional vectors. The first members are
unchanged because they are already orthonormal.
\end{proof}

\begin{example}[Gram--Schmidt in \(\mathbb R^2\)]
Starting with
\[
v_1=(1,1),
\qquad
v_2=(1,0),
\]
we obtain
\[
e_1=\frac1{\sqrt2}(1,1),
\]
and
\[
u_2=(1,0)-\left\langle(1,0),e_1\right\rangle e_1
=\left(\frac12,-\frac12\right).
\]
Thus
\[
e_2=\frac1{\sqrt2}(1,-1).
\]
The original and orthonormal lists span the same plane, but the latter makes
lengths, coordinates, and projections immediate.
\end{example}

\section{Orthogonal Complements and Projections}

An arbitrary direct-sum complement of a subspace is seldom canonical. The
orthogonal complement is canonical because it is determined by the inner
product alone; it separates a vector into a component in the subspace and a
perpendicular error.

\begin{proposition}[Orthogonal complements]
If \(W\leq V\), then \(W^\perp\) is a subspace. In finite dimensions,
\[
V=W\oplus W^\perp,
\qquad
\dim W+\dim W^\perp=\dim V,
\qquad
(W^\perp)^\perp=W.
\]
\end{proposition}
\begin{proof}
Closure of \(W^\perp\) under linear combinations follows from linearity in
the first variable. Choose an orthonormal basis
\[
(e_1,\ldots,e_r)
\]
of \(W\). For \(v\in V\), define
\[
P_Wv=\sum_{i=1}^r\langle v,e_i\rangle e_i.
\]
Then \(P_Wv\in W\), while, for each \(i\),
\[
\langle v-P_Wv,e_i\rangle=0.
\]
Thus \(v-P_Wv\in W^\perp\), proving \(V=W+W^\perp\). A vector in the
intersection is orthogonal to itself and is therefore zero, so the sum is
direct. Taking dimensions gives the formula. Since \(W\subseteq
(W^\perp)^\perp\) and the two spaces have the same dimension, they are
equal.
\end{proof}

\begin{definition}
The linear operator
\[
P_W:V\longrightarrow V,
\qquad
P_Wv=\sum_{i=1}^r\langle v,e_i\rangle e_i,
\]
where \(e_1,\ldots,e_r\) is any orthonormal basis of \(W\), is the
\emph{orthogonal projection onto \(W\)}; its image is \(W\).
\end{definition}

The preceding proof shows that the formula depends only on the unique
decomposition
\[
v=P_Wv+(v-P_Wv),
\qquad
P_Wv\in W,
\quad v-P_Wv\in W^\perp,
\]
and therefore not on the chosen orthonormal basis.

\begin{theorem}[Best approximation]
For \(v\in V\) and \(w\in W\),
\[
\lVert v-w\rVert^2
=\lVert v-P_Wv\rVert^2+\lVert P_Wv-w\rVert^2.
\]
Consequently \(P_Wv\) is the unique point of \(W\) nearest to \(v\).
\end{theorem}
\begin{proof}
The vector \(v-P_Wv\) lies in \(W^\perp\), whereas \(P_Wv-w\) lies in
\(W\). They are orthogonal and their sum is \(v-w\), so Pythagoras gives the
formula. The second summand is nonnegative and vanishes exactly when
\(w=P_Wv\), proving uniqueness.
\end{proof}

\begin{example}[Projection onto a line]
Let
\[
W=\operatorname{span}(1,1,1)\subseteq\mathbb R^3.
\]
The unit vector \(e=(1,1,1)/\sqrt3\) gives
\[
P_W(1,2,3)=\langle(1,2,3),e\rangle e=2(1,1,1).
\]
The error
\[
(-1,0,1)
\]
has dot product zero with \((1,1,1)\), visibly confirming the orthogonal
decomposition.
\end{example}

\section{Linear Functionals and Adjoints}

The algebraic dual \(V^\vee\), developed in Section~6.15, consists of all
linear measurements on \(V\). In finite-dimensional inner-product spaces,
the Riesz theorem says that every such measurement is obtained by taking an
inner product with one uniquely determined vector. It is this theorem that
makes the adjoint an intrinsic operator rather than merely a matrix operation.

\begin{theorem}[Finite-dimensional Riesz representation]
For every linear functional
\[
\varphi:V\longrightarrow F,
\]
there is a unique \(u\in V\) such that
\[
\varphi(v)=\langle v,u\rangle
\qquad(v\in V).
\]
\end{theorem}
\begin{proof}
Choose an orthonormal basis \(e_1,\ldots,e_n\). Put
\[
u=\sum_{j=1}^n\overline{\varphi(e_j)}e_j.
\]
If
\[
v=\sum_{j=1}^nx_je_j,
\]
then
\[
\langle v,u\rangle
=\sum_{j=1}^nx_j\varphi(e_j)
=\varphi(v).
\]
If both \(u\) and \(u'\) represent \(\varphi\), then
\[
\langle v,u-u'\rangle=0
\]
for every \(v\). Taking \(v=u-u'\) proves \(u=u'\).
\end{proof}

\begin{theorem}[Adjoint]
For every linear map
\[
T:V\longrightarrow W,
\]
there is a unique linear map
\[
T^*:W\longrightarrow V
\]
satisfying
\[
\langle Tv,w\rangle=\langle v,T^*w\rangle
\qquad(v\in V,\ w\in W).
\]
\end{theorem}
\begin{proof}
For fixed \(w\in W\), the function
\[
v\longmapsto\langle Tv,w\rangle
\]
is a linear functional on \(V\). Riesz representation supplies a unique
vector, denoted \(T^*w\), satisfying the displayed equality. For
\(a,b\in F\) and \(w_1,w_2\in W\), conjugate-linearity in the second
variable gives
\[
\langle Tv,aw_1+bw_2\rangle
=\overline a\langle Tv,w_1\rangle
+\overline b\langle Tv,w_2\rangle
=\langle v,aT^*w_1+bT^*w_2\rangle.
\]
Uniqueness in Riesz representation proves the linearity of \(T^*\). The
same uniqueness proves uniqueness of the adjoint.
\end{proof}

\noindent When \(V=W\), this is precisely the adjoint of the linear operator
\(T:V\to V\) considered below.  The point of allowing two spaces here is
that the adjoint reverses the direction of a map:
\[
T:V\longrightarrow W,
\qquad
T^*:W\longrightarrow V.
\]

\begin{proposition}[Adjoint identities and matrices]
The familiar adjoint identities remain valid for maps between possibly
different inner-product spaces, provided all sums and compositions are
defined.  Thus, if \(S,T:V\to W\) and \(a\in F\), then
\[
(S+T)^*=S^*+T^*,
\qquad
(aT)^*=\overline a\,T^*,
\qquad
(T^*)^*=T.
\]
If \(T:V\to W\) and \(S:W\to U\), then \(ST:V\to U\) and
\[
(ST)^*=T^*S^*:U\to V.
\]
If \(T:V\to W\) is an isomorphism, then
\[
(T^{-1})^*=(T^*)^{-1}.
\]
Here \(T^{-1}:W\to V\), so \((T^{-1})^*:V\to W\).  Likewise,
\(T^*:W\to V\), so \((T^*)^{-1}:V\to W\).
Relative to orthonormal bases, the matrix of \(T^*\) is
\[
[T^*]=[T]^*=\overline{[T]}^{\,t}.
\]
\end{proposition}
\begin{proof}
For the sum identity, the defining equation and linearity in the first
variable give
\[
\langle(S+T)v,w\rangle
=\langle v,(S^*+T^*)w\rangle.
\]
For the scalar identity, our convention gives
\[
\langle aTv,w\rangle
=a\langle Tv,w\rangle
=a\langle v,T^*w\rangle
=\langle v,\overline a\,T^*w\rangle.
\]
Over \(\mathbb R\), \(\overline a=a\).  For composition, if
\(T:V\to W\) and \(S:W\to U\), then
\[
\langle STv,w\rangle
=\langle Tv,S^*w\rangle
=\langle v,T^*S^*w\rangle.
\]
Finally,
\[
\langle T^*w,v\rangle_V
=\overline{\langle v,T^*w\rangle_V}
=\overline{\langle Tv,w\rangle_W}
=\langle w,Tv\rangle_W,
\]
so uniqueness shows that \((T^*)^*=T\).  For an isomorphism \(T:V\to W\),
\[
(T^{-1}T)^*=T^*(T^{-1})^*=I_V
\]
and
\[
(TT^{-1})^*=(T^{-1})^*T^*=I_W.
\]
Thus \((T^{-1})^*\) is the inverse of \(T^*\).
For orthonormal coordinate columns \(x,y\) and a matrix \(A=[T]\),
\[
\langle Ax,y\rangle
=x^tA^t\overline y
=\langle x,\overline A^{\,t}y\rangle.
\]
Thus the matrix of the adjoint is the conjugate transpose. Over
\(\mathbb R\), this is the ordinary transpose.
\end{proof}

\begin{definition}[Hermitian and symmetric matrices]
For \(A\in M_n(\mathbb C)\), its \emph{conjugate transpose} is
\[
A^*=\overline A^{\,\mathsf T}.
\]
The matrix \(A\) is \emph{Hermitian} if
\[
A^*=A.
\]
Equivalently, its entries satisfy
\[
a_{ij}=\overline{a_{ji}}
\qquad(1\leq i,j\leq n).
\]
In particular, every diagonal entry of a Hermitian matrix is real.  For
example,
\[
\begin{pmatrix}2&i\\-i&3\end{pmatrix}
\]
is Hermitian.  Over \(\mathbb R\), conjugation has no effect, and the
corresponding condition \(A^{\mathsf T}=A\) defines a \emph{symmetric}
matrix.
\end{definition}

\begin{remark}[Self-adjoint operators and matrices]
If \(\beta\) is an orthonormal basis of a finite-dimensional inner-product
space, then
\[
[T^*]_\beta=[T]_\beta^*.
\]
Consequently, over \(\mathbb C\),
\[
T\text{ is self-adjoint}
\quad\Longleftrightarrow\quad
[T]_\beta\text{ is Hermitian};
\]
over \(\mathbb R\), the corresponding statement uses symmetric matrices.
The orthonormality of \(\beta\) matters: in an arbitrary basis, the matrix
of a self-adjoint operator need not literally be Hermitian or symmetric.
\end{remark}

\begin{theorem}[Kernels, images, and adjoints]
For \(T:V\to W\),
\[
\ker T^*=(\operatorname{im}T)^\perp,
\qquad
\operatorname{im}T^*=(\ker T)^\perp.
\]
Consequently,
\[
\operatorname{rank}T=\operatorname{rank}T^*,
\]
and the equation \(Tv=w\) is solvable if and only if
\[
w\perp\ker T^*.
\]
\end{theorem}
\begin{proof}
For \(w\in W\),
\[
\begin{aligned}
w\in\ker T^*
&\Longleftrightarrow
\langle v,T^*w\rangle=0
&&\text{for every }v\in V \\
&\Longleftrightarrow
\langle Tv,w\rangle=0
&&\text{for every }v\in V.
\end{aligned}
\]
This is exactly \(w\in(\operatorname{im}T)^\perp\), proving the first
identity. Applying it to \(T^*\), using \((T^*)^*=T\) and taking a second
orthogonal complement in finite dimensions gives
\[
\operatorname{im}T^*=(\ker T)^\perp.
\]
The rank equality follows either by dimensions of these complements or by
rank--nullity. Finally,
\[
w\in\operatorname{im}T
\Longleftrightarrow
w\in(\ker T^*)^\perp,
\]
which is the asserted solvability criterion.
\end{proof}

\begin{proposition}[Characterization of orthogonal projections]
For a linear operator \(P:V\to V\), the following are equivalent:
\begin{enumerate}
\item \(P=P_W\) for some subspace \(W\);
\item \(P^2=P\) and \(P^*=P\).
\end{enumerate}
In this situation,
\[
W=\operatorname{im}P,
\qquad
\ker P=W^\perp.
\]
\end{proposition}
\begin{proof}
The formula for \(P_W\) shows directly that \(P_W^2=P_W\). If
\(v,w\in V\), expansion in an orthonormal basis of \(W\) gives
\[
\langle P_Wv,w\rangle=\langle v,P_Ww\rangle,
\]
so \(P_W^*=P_W\).

Conversely, suppose \(P^2=P=P^*\). The idempotence gives
\[
V=\operatorname{im}P\oplus\ker P.
\]
If \(x=Pu\in\operatorname{im}P\) and \(y\in\ker P\), then
\[
\langle x,y\rangle=\langle Pu,y\rangle
=\langle u,Py\rangle=0.
\]
Thus \(\ker P\subseteq(\operatorname{im}P)^\perp\). The direct-sum
dimension formula makes the two subspaces equal, so \(P\) is the orthogonal
projection onto its image.
\end{proof}

An ordinary algebraic projection only requires \(P^2=P\). It need not be
self-adjoint, and its kernel need not be perpendicular to its image. The
self-adjoint condition is exactly what distinguishes orthogonal projection.

\section{Self-Adjoint, Normal, Unitary, and Orthogonal Operators}

Adjoints identify the maps compatible with inner-product geometry. These
operator classes are important not merely for their names: normality is the
precise complex hypothesis for orthonormal diagonalization.

\begin{definition}
For an endomorphism \(T:V\to V\), we say that \(T\) is \emph{self-adjoint}
if
\[
T=T^*,
\]
\emph{skew-adjoint} if
\[
T^*=-T,
\]
and \emph{normal} if
\[
TT^*=T^*T.
\]
It is \emph{positive} if it is self-adjoint and
\[
\langle Tv,v\rangle\geq0
\qquad(v\in V),
\]
and \emph{positive-definite} if equality occurs only for \(v=0\). A complex
operator is \emph{unitary}, and a real operator is \emph{orthogonal}, if
\[
T^*T=I.
\]
\end{definition}

Self-adjoint matrices are real symmetric or complex Hermitian. The terms
unitary and orthogonal distinguish the base field because the latter is the
real specialization of the former.

\begin{proposition}[Isometries and unitary operators]
For an endomorphism \(T\) of a finite-dimensional inner-product space, the
following are equivalent:
\begin{enumerate}
\item \(T\) is unitary or orthogonal, as appropriate to the field;
\item \(T^*T=I\);
\item
\[
\langle Tu,Tv\rangle=\langle u,v\rangle
\qquad(u,v\in V);
\]
\item
\[
\lVert Tv\rVert=\lVert v\rVert
\qquad(v\in V).
\]
\end{enumerate}
In this case \(T\) is invertible and
\[
T^{-1}=T^*.
\]
Relative to an orthonormal basis, its matrix \(A\) satisfies
\[
A^*A=AA^*=I.
\]
Thus the columns and rows of \(A\) are orthonormal bases.
\end{proposition}
\begin{proof}
The equivalence of \(1\) and \(2\) is the definition. Moreover,
\[
\langle Tu,Tv\rangle=\langle u,T^*Tv\rangle,
\]
so \(2\) implies \(3\), and \(3\) implies \(4\). Conversely, the polarization
identity shows that \(4\) determines the same inner products, so \(4\) implies
\(3\). If \(3\) holds, then
\[
\langle u,(T^*T-I)v\rangle=0
\]
for every \(u,v\); taking \(u=(T^*T-I)v\) proves \(2\).

Condition \(4\) makes \(T\) injective. Finite dimensionality then makes it
surjective, so it is invertible; from \(T^*T=I\) we obtain \(T^*=T^{-1}\),
and hence \(TT^*=I\). The matrix identities follow in orthonormal bases.
Their entries state exactly that both the columns and rows are orthonormal.
\end{proof}

\begin{proposition}[Basic consequences of the adjoint]
Every operator \(T^*T\) is self-adjoint and positive. If \(T\) is normal,
then
\[
\lVert Tv\rVert=\lVert T^*v\rVert
\qquad(v\in V),
\]
and
\[
\ker T=\ker T^*.
\]
Self-adjoint and unitary or orthogonal operators are normal.
\end{proposition}
\begin{proof}
The adjoint identities give
\[
(T^*T)^*=T^*T,
\]
and
\[
\langle T^*Tv,v\rangle=\langle Tv,Tv\rangle\geq0.
\]
If \(T\) is normal, then
\[
\lVert Tv\rVert^2
=\langle T^*Tv,v\rangle
=\langle TT^*v,v\rangle
=\lVert T^*v\rVert^2.
\]
This equality proves the kernel identity. Self-adjointness makes the two
products identical, while for a unitary or orthogonal operator they both
equal \(I\), proving the final assertion.
\end{proof}

\begin{proposition}[Eigenvalue geometry]
Let \(T\) be an endomorphism of a finite-dimensional inner-product space.
\begin{enumerate}
\item If \(T\) is self-adjoint, every eigenvalue of \(T\) is real.
\item If \(T\) is unitary or orthogonal, every eigenvalue \(\lambda\) has
\[
|\lambda|=1.
\]
\item If \(T\) is normal on a complex inner-product space and
\[
Tv=\lambda v,
\]
then
\[
T^*v=\overline\lambda v.
\]
Consequently, eigenvectors of a normal operator belonging to distinct
eigenvalues are orthogonal.
\end{enumerate}
\end{proposition}
\begin{proof}
For a self-adjoint \(T\) and a nonzero eigenvector \(v\),
\[
\lambda\langle v,v\rangle
=\langle Tv,v\rangle
=\langle v,Tv\rangle
=\overline\lambda\langle v,v\rangle,
\]
so \(\lambda=\overline\lambda\). If \(T\) is unitary or orthogonal,
\[
\lVert v\rVert=\lVert Tv\rVert=|\lambda|\lVert v\rVert,
\]
proving the second claim. For normal \(T\), the preceding proposition gives
\[
\lVert(T-\lambda I)v\rVert
=\lVert(T^*-\overline\lambda I)v\rVert.
\]
The left side is zero, hence the third assertion follows. If also
\[
Tw=\mu w,
\]
then
\[
\lambda\langle v,w\rangle
=\langle Tv,w\rangle
=\langle v,T^*w\rangle
=\langle v,\overline\mu w\rangle
=\mu\langle v,w\rangle.
\]
Thus \(\langle v,w\rangle=0\) when \(\lambda\ne\mu\).
\end{proof}

For a normal operator on a complex inner-product space, the last assertion
may be recorded as
\[
\lambda\ne\mu
\quad\Longrightarrow\quad
E_\lambda\perp E_\mu.
\]
Here the proof above applies to arbitrary vectors in the two eigenspaces,
not only to a particular choice of eigenvectors.

\begin{example}[Why real eigenvalues are not enough]
On \(\mathbb C^2\) with its standard inner product, consider
\[
A=\begin{pmatrix}1&1\\0&2\end{pmatrix}.
\]
Its eigenvalues are \(1\) and \(2\), with eigenvectors
\[
v_1=(1,0),
\qquad
v_2=(1,1),
\]
respectively.  Thus \(A\) is diagonalizable, and all of its eigenvalues are
real.  Nevertheless,
\[
A^*=\begin{pmatrix}1&0\\1&2\end{pmatrix}\ne A,
\qquad
A^*A=\begin{pmatrix}1&1\\1&5\end{pmatrix}
\ne
\begin{pmatrix}2&2\\2&4\end{pmatrix}=AA^*,
\]
so \(A\) is neither self-adjoint nor normal.  Also,
\(\langle v_1,v_2\rangle=1\), showing that eigenvectors for distinct
eigenvalues of an arbitrary diagonalizable operator need not be orthogonal.

This example rules out both false implications
\[
\begin{gathered}
\text{real eigenvalues}\ \not\Longrightarrow\ \text{self-adjoint},\\
\text{diagonalizable with real eigenvalues}\ \not\Longrightarrow\ \text{self-adjoint}.
\end{gathered}
\]
It also illustrates the danger of applying Gram--Schmidt globally to an
eigenbasis.  Starting with \((v_1,v_2)\), the process gives
\[
e_1=(1,0),
\qquad
e_2=v_2-\langle v_2,e_1\rangle e_1=(0,1).
\]
But
\[
Ae_2=(1,2)
\]
is not a scalar multiple of \(e_2\), so the resulting orthonormal basis is
not an eigenbasis.
\end{example}

\begin{remark}[Local, not global, Gram--Schmidt]
If
\[
V=\bigoplus_\lambda E_\lambda
\]
is an eigenspace decomposition, Gram--Schmidt can safely be applied inside
each individual eigenspace: linear combinations of vectors in
\(E_\lambda\) remain eigenvectors with eigenvalue \(\lambda\).  Its union
is an orthonormal basis of \(V\) precisely when the distinct eigenspaces are
mutually orthogonal.  The example shows why applying the process across
nonorthogonal eigenspaces can fail.  For a normal complex operator, the
orthogonality just proved makes the local procedure work globally; if an
eigenbasis is already orthonormal, the preceding corollary says there is
nothing for Gram--Schmidt to change.
\end{remark}

\begin{remark}[The real and complex cases]
A self-adjoint operator on a finite-dimensional real inner-product space has
real eigenvalues and an orthonormal eigenbasis, as the real spectral theorem
below states.  A real normal operator need not be diagonalizable over
\(\mathbb R\).  Indeed,
\[
R_\theta=
\begin{pmatrix}
\cos\theta&-\sin\theta\\
\sin\theta&\cos\theta
\end{pmatrix}
\]
is orthogonal and hence normal, but for
\(\theta\notin\{0,\pi\}\) it has no real eigenvalues.  The real analogue
of unitary diagonalization for general normal operators is orthogonal block
diagonalization, with real \(1\times1\) blocks and \(2\times2\)
rotation-type blocks.  Thus the full normal spectral theorem is naturally a
complex statement.
\end{remark}

\section{Spectral Theorems}

The spectral theorem identifies exactly when ordinary diagonalization can be
made geometric. Its finite-dimensional proof requires no Hilbert-space
theory, but the complex result uses the Fundamental Theorem of Algebra to
produce an eigenvalue.

\begin{remark}[The Fundamental Theorem of Algebra]
The Fundamental Theorem of Algebra says that every nonconstant polynomial in
\(\mathbb C[x]\) has a root in \(\mathbb C\). We use it only to ensure that
the characteristic polynomial of a complex endomorphism has a root. A
Galois-theoretic proof appears later in Chapter~11; no complex analysis is
needed elsewhere in this chapter.
\end{remark}

\begin{theorem}[Complex spectral theorem]
For an endomorphism \(T\) of a finite-dimensional complex inner-product
space, the following are equivalent:
\begin{enumerate}
\item \(T\) is normal;
\item \(V\) has an orthonormal basis of eigenvectors of \(T\);
\item the matrix of \(T\) in some orthonormal basis is diagonal.
\end{enumerate}
Equivalently, \(T\) is unitarily diagonalizable.
\end{theorem}
\begin{proof}
The equivalence of \(2\) and \(3\) is the definition of a matrix in a chosen
basis. A diagonal matrix commutes with its conjugate transpose, so \(2\)
implies \(1\). Conversely, suppose \(T\) is normal and argue by induction on
\(\dim V\). The case \(\dim V=0\) is clear. The Fundamental Theorem of
Algebra supplies an eigenvalue \(\lambda\), and let
\[
E=\ker(T-\lambda I).
\]
If \(z\in E^\perp\) and \(e\in E\), then
\[
\langle Tz,e\rangle
=\langle z,T^*e\rangle
=\langle z,\overline\lambda e\rangle=0.
\]
Thus \(E^\perp\) is \(T\)-invariant. The same calculation with \(T^*\)
shows that it is \(T^*\)-invariant, so the restriction of \(T\) to
\(E^\perp\) remains normal. It has smaller dimension unless \(E=V\), and
the induction hypothesis supplies an orthonormal eigenbasis of \(E^\perp\).
An orthonormal basis of \(E\), obtained by Gram--Schmidt, together with that
basis gives an orthonormal eigenbasis of \(V\).
\end{proof}

\begin{corollary}[Real spectral theorem]
For an endomorphism \(T\) of a finite-dimensional real inner-product space,
the following are equivalent:
\begin{enumerate}
\item \(T\) is self-adjoint;
\item \(V\) has an orthonormal basis of real eigenvectors of \(T\);
\item \(T\) is orthogonally diagonalizable.
\end{enumerate}
\end{corollary}
\begin{proof}
The last two conditions are equivalent by coordinates, and either one plainly
implies self-adjointness. Conversely, complexify \(V\) and extend the real
inner product to the usual Hermitian inner product on \(V_\mathbb C\). The
complexified operator is self-adjoint, hence normal; the complex spectral
theorem applies. Its eigenvalues are real by the preceding eigenvalue
proposition. If
\[
T_\mathbb C(x+iy)=\lambda(x+iy),
\]
then
\[
Tx=\lambda x,
\qquad
Ty=\lambda y.
\]
Thus the complex eigenspaces are complexifications of real eigenspaces.
They span \(V_\mathbb C\), so the real eigenspaces span \(V\). Distinct
eigenspaces are orthogonal, and Gram--Schmidt within each one produces a
real orthonormal eigenbasis.
\end{proof}

\begin{remark}[Spectral-theorem implications]
For a finite-dimensional complex inner-product space, the main implications
are
\[
\text{self-adjoint}
\ \Longrightarrow\
\text{normal}
\ \Longrightarrow\
\text{unitarily diagonalizable}
\ \Longrightarrow\
\text{diagonalizable}.
\]
The reverse implications fail without additional hypotheses.  More exactly,
\[
T\text{ is normal}
\quad\Longleftrightarrow\quad
T\text{ has an orthonormal eigenbasis},
\]
and
\[
\begin{gathered}
T\text{ is self-adjoint}\\
\Longleftrightarrow\\
T\text{ has an orthonormal eigenbasis with real eigenvalues}.
\end{gathered}
\]
The latter equivalence explains the geometric content that is missing from
ordinary diagonalizability.  It also gives the useful converse
\[
T\text{ normal with real spectrum}
\quad\Longrightarrow\quad
T\text{ self-adjoint}:
\]
in an orthonormal eigenbasis, its diagonal matrix has real entries and hence
equals its conjugate transpose.  At the matrix level, this says
\[
\begin{gathered}
A\text{ is Hermitian}\\
\Longleftrightarrow\\
A\text{ is unitarily diagonalizable with real eigenvalues}.
\end{gathered}
\]
\end{remark}

\begin{example}[Normal diagonalization]
The matrix
\[
\begin{pmatrix}2&0\\0&i\end{pmatrix}
\]
is normal and already unitarily diagonal. By contrast,
\[
\begin{pmatrix}1&a\\0&1\end{pmatrix}
\]
is diagonalizable only when \(a=0\). For \(a\ne0\), it does not commute
with its conjugate transpose and is not normal. The spectral theorem
therefore excludes precisely the Jordan-block obstruction described in
Chapter~6.
\end{example}

\section{Singular Value Decomposition}

Eigenvalue diagonalization applies only to endomorphisms.  Singular value
decomposition applies to a map between two possibly different spaces.  Its
conceptual path is
\[
\begin{gathered}
\text{adjoint}\longrightarrow T^*T
\longrightarrow\text{positive self-adjoint operator}\\
\longrightarrow\text{spectral theorem}\longrightarrow\text{SVD}.
\end{gathered}
\]
Here the adjoint constructed in the preceding section has the defining
property
\[
\langle Tv,w\rangle_W=\langle v,T^*w\rangle_V
\qquad(v\in V,\ w\in W).
\]
Thus, although \(T\) need not be an endomorphism, the two composites are
endomorphisms of different spaces:
\[
T^*T:V\longrightarrow V,
\qquad
TT^*:W\longrightarrow W.
\]

\begin{proposition}[The positive operators associated to a map]
For a linear map \(T:V\to W\), both \(T^*T\) and \(TT^*\) are self-adjoint
and positive.  In particular, all their eigenvalues are real and
nonnegative.
\end{proposition}
\begin{proof}
The generalized adjoint identities just proved give
\[
(T^*T)^*=T^*(T^*)^*=T^*T,
\qquad
(TT^*)^*=(T^*)^*T^*=TT^*.
\]
Moreover, for \(v\in V\) and \(w\in W\),
\[
\langle T^*Tv,v\rangle_V=\langle Tv,Tv\rangle_W
=\lVert Tv\rVert^2\geq0,
\]
and
\[
\langle TT^*w,w\rangle_W=\langle T^*w,T^*w\rangle_V
=\lVert T^*w\rVert^2\geq0.
\]
If, for example, \(T^*Tv=\lambda v\) with \(v\ne0\), then
\[
\lambda\lVert v\rVert^2=\langle T^*Tv,v\rangle_V
=\lVert Tv\rVert^2.
\]
The left side is real because \(T^*T\) is self-adjoint (or by the
eigenvalue geometry proposition), and the right side is nonnegative.
Hence \(\lambda\) is real and nonnegative.  The same argument applies to
\(TT^*\).
\end{proof}

\begin{proposition}[Nonzero spectra]
The operators \(T^*T\) and \(TT^*\) have the same nonzero eigenvalues, with
the same multiplicities.
\end{proposition}
\begin{proof}
Let \(\lambda>0\).  The map \(T\) sends the \(\lambda\)-eigenspace of
\(T^*T\) into the \(\lambda\)-eigenspace of \(TT^*\), since
\[
TT^*(Tv)=T(T^*Tv)=\lambda Tv.
\]
It is injective on this eigenspace: if \(Tv=0\), then
\(\lambda v=T^*Tv=0\).  Conversely, if \(TT^*w=\lambda w\), then
\[
T\left(\frac{T^*w}{\lambda}\right)=w.
\]
Moreover,
\[
T^*T\left(\frac{T^*w}{\lambda}\right)
=\frac{T^*(TT^*w)}{\lambda}
=\lambda\left(\frac{T^*w}{\lambda}\right),
\]
so this preimage belongs to the \(\lambda\)-eigenspace of \(T^*T\).
Thus \(T\) is an isomorphism between the two \(\lambda\)-eigenspaces.
Since both operators are self-adjoint and hence diagonalizable, the
dimensions of these eigenspaces are their multiplicities.
\end{proof}

\begin{definition}
List the eigenvalues of \(T^*T\), with multiplicity, as
\[
\lambda_1,\ldots,\lambda_n,
\qquad n=\dim V.
\]
The numbers
\[
\sigma_i=\sqrt{\lambda_i}
\]
are the \emph{singular values} of \(T\).  They are well-defined nonnegative
real numbers by the preceding proposition.  The nonzero singular values are
called the \emph{positive singular values}.
\end{definition}

The norm identity also records exactly where the zero singular values occur:
\[
\ker(T^*T)=\ker T.
\]
Indeed, \(Tv=0\) plainly implies \(T^*Tv=0\), while \(T^*Tv=0\) implies
\(\lVert Tv\rVert^2=0\), hence \(Tv=0\).  Consequently, the number of
positive singular values, counted with multiplicity, is
\[
\dim V-\dim\ker T=\operatorname{rank}T.
\]

\begin{theorem}[Singular value decomposition]
Let \(T:V\to W\) be a linear map of finite-dimensional real or complex
inner-product spaces.  Write
\[
n=\dim V,\qquad m=\dim W,\qquad r=\operatorname{rank}T.
\]
There are orthonormal bases
\[
(v_1,\ldots,v_n)\quad\text{of }V,
\qquad
(u_1,\ldots,u_m)\quad\text{of }W,
\]
and positive singular values \(\sigma_1,\ldots,\sigma_r\) such that
\[
Tv_i=\sigma_i u_i\quad(1\leq i\leq r),
\qquad
Tv_i=0\quad(r<i\leq n).
\]
Equivalently, relative to these bases the matrix of \(T\) is the rectangular
diagonal matrix \(\Sigma=(\Sigma_{ij})\) of size \(m\times n\), where
\[
\Sigma_{ij}=
\begin{cases}
\sigma_i,& i=j\leq r,\\
0,&\text{otherwise}.
\end{cases}
\]
\end{theorem}
\begin{proof}
By the real or complex spectral theorem, the positive self-adjoint operator
\(T^*T:V\to V\) has an orthonormal eigenbasis.  Arrange it so that
\[
T^*Tv_i=\sigma_i^2v_i\quad(1\leq i\leq r),
\]
where \(\sigma_i>0\), and so that \(v_{r+1},\ldots,v_n\) span the
zero-eigenspace.  The kernel identity above says that this zero-eigenspace
is \(\ker T\), so \(Tv_i=0\) for \(i>r\).

For \(1\leq i\leq r\), define
\[
u_i=\frac{Tv_i}{\sigma_i}\in W.
\]
These vectors are orthonormal.  Indeed, since the singular values are real,
\[
\begin{aligned}
\langle u_i,u_j\rangle_W
&=\frac{1}{\sigma_i\sigma_j}\langle Tv_i,Tv_j\rangle_W\\
&=\frac{1}{\sigma_i\sigma_j}\langle T^*Tv_i,v_j\rangle_V\\
&=\frac{\sigma_i}{\sigma_j}\langle v_i,v_j\rangle_V
=\delta_{ij}.
\end{aligned}
\]
Extend \(u_1,\ldots,u_r\) to an orthonormal basis
\(u_1,\ldots,u_m\) of \(W\).  The defining equation for the \(u_i\)'s and
the kernel calculation give the asserted formulas for \(Tv_i\), and these
formulas say exactly that the matrix is \(\Sigma\).
\end{proof}

For matrices, work in fixed orthonormal coordinate bases and write \(A=[T]\).
Let \(V\) and \(U\) be the matrices whose
columns are respectively the coordinate columns of \(v_1,\ldots,v_n\) and
\(u_1,\ldots,u_m\).  These matrices are unitary over \(\mathbb C\) and
orthogonal over \(\mathbb R\).  Since the matrix in the singular-vector
bases is \(\Sigma\),
\[
\Sigma=U^*AV,
\qquad\text{and hence}\qquad
A=U\Sigma V^*.
\]
Over \(\mathbb R\), this reads
\[
A=U\Sigma V^{\mathsf T}.
\]

\subsection{Spectral decomposition and SVD}

The case \(T:V\to V\) makes especially transparent both what SVD inherits
from the spectral theorem and what information it discards.

\begin{proposition}[Spectral decomposition]
Let \(T:V\to V\) be self-adjoint, let
\(\lambda_1,\ldots,\lambda_s\) be its distinct eigenvalues, and let
\(P_{\lambda_j}\) be the orthogonal projection onto
\(E_{\lambda_j}\).  Then
\[
T=\sum_{j=1}^s\lambda_jP_{\lambda_j},
\qquad
I=\sum_{j=1}^sP_{\lambda_j},
\]
and
\[
P_{\lambda_i}P_{\lambda_j}=0\quad(i\ne j),
\qquad
P_{\lambda_j}^*=P_{\lambda_j}.
\]
\end{proposition}
\begin{proof}
The real or complex spectral theorem gives the orthogonal direct sum
\[
V=\bigoplus_{j=1}^sE_{\lambda_j}.
\]
For \(v=\sum_jv_j\), with \(v_j\in E_{\lambda_j}\), the projections satisfy
\(P_{\lambda_j}v=v_j\).  Hence
\[
Tv=\sum_j\lambda_jv_j
=\sum_j\lambda_jP_{\lambda_j}v,
\]
and the remaining identities follow from the orthogonal direct sum and the
fact that orthogonal projections are self-adjoint.
\end{proof}

This formula is the coordinate-free form of unitary or orthogonal
diagonalization.  It also describes the SVD of a self-adjoint operator.  If
\(v_1,\ldots,v_n\) is an orthonormal eigenbasis with
\[
Tv_i=\lambda_iv_i,
\]
then \(T^*T=T^2\), so
\[
T^*Tv_i=\lambda_i^2v_i.
\]
Thus the corresponding singular values are
\[
\sigma_i=|\lambda_i|.
\]
For \(\lambda_i\ne0\), set
\[
u_i=\operatorname{sgn}(\lambda_i)v_i.
\]
Then \(u_i\) is a unit vector and
\[
Tv_i=\sigma_i u_i.
\]
When \(\lambda_i=0\), take \(u_i=v_i\); then again
\(Tv_i=0=\sigma_i u_i\).  The \(u_i\)'s form an orthonormal basis, so this
is an SVD.  The ordinary spectral decomposition retains the sign of every
real eigenvalue, whereas SVD replaces it by its absolute value.

If \(T\) is positive semidefinite as well as self-adjoint, then
\(\lambda_i\geq0\), and hence \(\sigma_i=\lambda_i\).  In this case one
may take \(u_i=v_i\): the SVD and spectral decomposition essentially
coincide.  At the matrix level, a Hermitian positive semidefinite matrix has
\[
A=U\Lambda U^*,
\qquad
\Lambda=\operatorname{diag}(\lambda_1,\ldots,\lambda_n),
\qquad
\lambda_i\geq0,
\]
and this is simultaneously a spectral decomposition and an SVD.

\begin{remark}[A brief polar-decomposition viewpoint]
For an endomorphism \(T\), the positive operator
\[
|T|=(T^*T)^{1/2}
\]
is defined by applying the spectral theorem to \(T^*T\).  Its eigenvalues
are the singular values of \(T\).  When \(T\) is self-adjoint, \(|T|\)
replaces each eigenvalue \(\lambda\) by \(|\lambda|\).  This is the
starting point for polar decomposition and another direct bridge between the
spectral theorem and SVD.
\end{remark}

\begin{example}[An SVD calculation]
Consider the real \(3\times2\) matrix
\[
A=\begin{pmatrix}1&1\\1&1\\0&0\end{pmatrix}.
\]
We first compute the operator on the domain:
\[
A^{\mathsf T}A
=\begin{pmatrix}1&1&0\\1&1&0\end{pmatrix}
\begin{pmatrix}1&1\\1&1\\0&0\end{pmatrix}
=\begin{pmatrix}2&2\\2&2\end{pmatrix}.
\]
Its characteristic polynomial is
\[
\det\begin{pmatrix}2-\lambda&2\\2&2-\lambda\end{pmatrix}
=\lambda(\lambda-4),
\]
so its eigenvalues are \(4\) and \(0\).  Solving the two eigenspace
equations gives
\[
(A^{\mathsf T}A-4I)\begin{pmatrix}x\\y\end{pmatrix}=0
\quad\Longleftrightarrow\quad x=y,
\qquad
A^{\mathsf T}A\begin{pmatrix}x\\y\end{pmatrix}=0
\quad\Longleftrightarrow\quad x=-y.
\]
Normalizing vectors in these two directions gives the orthonormal eigenbasis
\[
v_1=\frac1{\sqrt2}\begin{pmatrix}1\\1\end{pmatrix},
\qquad
v_2=\frac1{\sqrt2}\begin{pmatrix}1\\-1\end{pmatrix}.
\]
Thus the singular values are \(\sigma_1=\sqrt4=2\) and
\(\sigma_2=\sqrt0=0\).  The vectors \(v_1,v_2\) are the right singular
vectors; in particular, \(v_2\) spans \(\ker A\).

For the positive singular value, the corresponding left singular vector is
\[
u_1=\frac{Av_1}{2}
=\frac12\begin{pmatrix}\sqrt2\\\sqrt2\\0\end{pmatrix}
=\begin{pmatrix}1/\sqrt2\\1/\sqrt2\\0\end{pmatrix}.
\]
Complete it to an orthonormal basis of \(\mathbb R^3\), for example by
choosing
\[
u_2=\begin{pmatrix}1/\sqrt2\\-1/\sqrt2\\0\end{pmatrix},
\qquad
u_3=\begin{pmatrix}0\\0\\1\end{pmatrix}.
\]
Therefore
\[
U=\begin{pmatrix}
1/\sqrt2&1/\sqrt2&0\\
1/\sqrt2&-1/\sqrt2&0\\
0&0&1
\end{pmatrix},
\qquad
\Sigma=\begin{pmatrix}2&0\\0&0\\0&0\end{pmatrix},
\qquad
V=\begin{pmatrix}
1/\sqrt2&1/\sqrt2\\
1/\sqrt2&-1/\sqrt2
\end{pmatrix}.
\]
Both \(U\) and \(V\) are orthogonal.  Finally, direct multiplication gives
\[
U\Sigma V^{\mathsf T}
=2u_1v_1^{\mathsf T}
=2\begin{pmatrix}1/\sqrt2\\1/\sqrt2\\0\end{pmatrix}
\begin{pmatrix}1/\sqrt2&1/\sqrt2\end{pmatrix}
=\begin{pmatrix}1&1\\1&1\\0&0\end{pmatrix}=A.
\]
\end{example}

\begin{remark}[Why SVD matters]
The SVD is also a practical tool.  It underlies least-squares problems and
the Moore--Penrose pseudoinverse, best low-rank approximation, and principal
component analysis for dimensionality reduction.  Keeping only the largest
singular values can compress images and other data, while discarding small
ones is useful in denoising and signal processing.  Singular values also
measure numerical conditioning, so SVD is central in numerical linear
algebra and in many applications across computer science and engineering.
\end{remark}

\section*{Exercises}
\begin{exercise}[Basic] Verify that the standard formula on \(\mathbb C^n\)
is an inner product under the convention of this chapter.\end{exercise}
\begin{exercise}[Core] Use the complex polarization identity to recover the
standard inner product of two specified vectors in \(\mathbb C^2\) from their
norms.\end{exercise}
\begin{exercise}[Core] Apply Gram--Schmidt to
\[
(1,1,0),\ (1,0,1),\ (0,1,1)
\]
in \(\mathbb R^3\).\end{exercise}
\begin{exercise}[Core] Find the orthogonal projection of \((1,2,3)\) onto
the span of \((1,1,1)\), and verify its best-approximation property for an
arbitrary point on that line.\end{exercise}
\begin{exercise}[Core] Compute the adjoint of
\[
\begin{pmatrix}1&i\\0&2\end{pmatrix}
\]
on \(\mathbb C^2\), and determine its kernel.\end{exercise}
\begin{exercise}[Further] Prove that a positive-definite operator is
invertible.\end{exercise}
\begin{exercise}[Further, Challenging] Find an SVD of
\[
\begin{pmatrix}1&1\\0&0\end{pmatrix}
\]
and identify its nonzero singular subspaces.\end{exercise}
